\section{Equivariant \texorpdfstring{$K$}{K}-theory and fixed-points}\label{section: equivariant K-theory and fixed points}

We now set conventions regarding equivariant topological $K$-theory. Here, we recall both the usual topological approach through homotopy theory \S \ref{section: generalized cohomology theories in topology}, as well as the algebraic approach via localizing invariants \S \ref{section: algebraic cohomology theories as localizing invariants}. We also take this as an opportunity to recall how these approaches relate -- they are equivalent on finite type schemes over $k=\C$. We emphasize that we work with genuine equivariance. 

After recalling our notion of fixed-point schemes in \S \ref{section: fixed-point schemes}, we finally present in \S \ref{section: K-theory and HP and cohomology of fixed points} our reinterpretation of classical fixed-point localization in terms of the trace to equivariant periodic cyclic homology. The main statement is Theorem \ref{theorem: fibers of topological K-theory, HP and fixed-points}. We find this formulation both conceptually clear and useful in practice. In particular, the result holds without restrictions for any (possibly singular) $X$ with respect to an arbitrary action of $T$. This formulation further gives the localization map a clear global meaning.






\subsection{Generalized cohomology theories in topology}\label{section: generalized cohomology theories in topology}



\subsubsection{Generalized (co)homology theories.}
Recall the notion of a generalized (co)homology theory in topology \cite{May96}. These are representable by spectra. Multiplicative generalized cohomology theories are represented by ring spectra. Given such spectrum $E$, one may evaluate it on any topological space $X$ by taking the mapping space $\map(E, -)$ resp. $\map(-, E)$ and collecting its homotopy groups into a $\Z$-graded abelian group:
\begin{equation*}
    E_{\bullet}(X) := \pi_{\bullet}(\map(E, X)) \qquad \text{resp.} \qquad E^{\bullet}(X) := \pi_{\bullet}(\map(X, E)).
\end{equation*}
These respectively give the $E$-homology resp $E$-cohomology of $X$.
Our exposition is biased towards the contravariant cohomological picture, which carries a natural ring structure when $E$ is multiplicative. We focus on this case from now on.


\begin{remark}\label{remark: sum and product convention for generalized cohomology theories}
To regard the above $E^{\bullet}(X)$ as actual ring, one needs to choose how to interpret the grading direction: for example, it is a matter of convention whether to take $\bigoplus_{\bullet \in \Z} \pi_{\bullet}(-)$ or $\prod_{\bullet \in \Z} \pi_{\bullet}(-)$.
%
In many situations -- such as ordinary cohomology of finite CW complexes -- this distinction is vacuous. However, in infinite-dimensional contexts and more general cohomology theories, one needs to treat this seriously. The correct cohomological convention is to take $\prod_{\bullet \in \Z}$. 
\end{remark}


\begin{notation}
Here and below, we use the following convention. For any coefficient ring $\k$, we write $E(-; \k)$ for the version with $\k$-coefficients, represented by the smash product of $E$ with the Moore space of $\k$.  When no coefficients are specified, we have $E(-) = E(-; \Z)$.    
\end{notation}



Given a Lie group $G$, one can axiomatically consider multiplicative $G$-equivariant generalized cohomology theories \cite[\S XIII]{May96}. These are representable by $G$-equivariant ring spectra.





\subsubsection{Generalized cohomology of colimits.}\label{subsubsection: K-theory of colimits}
Recall the following principle \cite[\S 19.4]{May99}, valid for any generalized cohomology theory.
\begin{recollection}
Let $j \in \Z$ and let $X$ be a topological space, given as an $\N_0$-indexed colimit
\begin{equation*}
    X = \colim_{i} X_i.
\end{equation*}
Then there is a short exact sequence
\begin{equation*}
0 \to {\lim_i}^1 (E^{j-1}(X_i)) \to E^j(X) \to \lim_i E^j(X_i) \to 0.    
\end{equation*}
Moreover, if the system $(E^{j-1}(X_i))_{i}$ satisfies the Mittag--Leffler condition, the ${\lim}^1$-term vanishes and we deduce
\begin{equation*}
    E^j(X) \xrightarrow{\cong} \lim_i E^j(X_i).
\end{equation*}
\end{recollection}

\begin{remark}\label{remark: Mittag--Leffler and limit of rings}
Under the Mittag--Leffler assumption and using the correct (product) convention from Remark \ref{remark: sum and product convention for generalized cohomology theories}, we deduce a ring isomorphism
\begin{equation*}
    E^{\bullet}(X) \cong \lim_i E^{\bullet}(X_i).
\end{equation*}
\end{remark}

\begin{remark}\label{remark: topological cellularity}
In particular, assume that $X$ is a CW complex with even cells only, together with its filtration by skeleta $\sk_i X$. Then, for any fixed $j$, all the maps in the inverse system $(E^{j}(X_i), i \in \N_0)$ are surjective, and the Mittag--Leffler condition is trivially satisfied.
\end{remark}

\begin{remark}
For a scheme $X$ over $\C$, consider its complex points $X(\C)$ with analytic topology. Given a generalized cohomology theory $E$, we then denote $E(X) := E(X(\C))$. We use similar convention in the equivariant case.

Assume that $X$ is a proper scheme over $\C$ that possesses a paving by affine spaces. Then its underlying topological space $X(\C)$ a CW complex with even cells only, and Remark \ref{remark: topological cellularity} applies.    
\end{remark}





\subsubsection{Cohomology.}
In particular, we denote by $H\Z$ the Eilenberg--Mac-Lane spectrum representing singular cohomology, and $H\k$ the version with coefficients. It represents singular cohomology; for a topological space $X$, we denote it $H(X) := H_{\sing}(X) := \map(X, H\Z)$. Similarly with other coefficients.

\begin{convention}
Given a topological space $X$, we write
\begin{equation*}
    \widehat{H}^{\bullet}(X) := \prod_{i \geq 0} H^{i}(X).
\end{equation*}
Here, the convention is that we take the direct product of all cohomology groups rather then their direct sum; see Remarks \ref{remark: sum and product convention for generalized cohomology theories}, \ref{remark: Mittag--Leffler and limit of rings}. 
%
It would be for the best to write simply $H^{\bullet}$ without any adornments; however, since the opposite convention is quite pervasive in geometric representation theory, we write $\widehat{H}^{\bullet}$ to avoid confusion.
\end{convention}


Cohomology extends to the equivariant setup through the Borel construction. Given a topological group $G$, denote $EG \to BG$ the classifying space of $G$ with its universal $G$-bundle. For any $G$-space $X$, its equivariant cohomology is defined by $H_G(X) := H(EG \times^G X)$. 






\subsubsection{Topological $K$-theory.}
We denote by $BU$ the complex topological $K$-theory spectrum, i.e. the ring spectrum representing topological $K$-theory. Given a topological space $X$, we denote its topological $K$-theory
\begin{equation*}
    K^{\top}(X) := KU(X) := \map(X, BU).
\end{equation*}
We write $K^{\top, \bullet}(X) := \pi_{\bullet}(K^{\top}(X))$ with $\bullet \in \Z$ for its homotopy groups, which form a graded-commutative ring.

Topological $K$-theory genuinely extends to the equivariant context. 
\begin{recollection}[Genuine equivariant topological $K$-theory]\label{recollection: genuine topological K-theory}
Let $G$ be a compact topological group. The genuine equivariant $K$-theory was firstly defined by \cite{Seg68}. These ideas were further developed in terms of the representing spectrum $BU_G$, see \cite[\S XIV, in particular \S XIV.3 and \S XIV.4]{May96}. Given any $G$-space $X$, we thus have equivariant topological $K$-theory
\begin{equation*}
    K^{\top}_G(X) := KU_G(X) := \map_G(X, BU_G).
\end{equation*}
We will use this genuine definition. The discussion of previous paragraphs works in this context. The definition extends to the case when the group $G$ is a complex Lie group by passing to its maximal compact subgroup $M$; see also \cite[\S 2.1.2]{HLP20}. For finite $G$-CW complexes, the above can be also defined directly from the category of equivariant topological vector bundles on $X$ as in \cite{Weib13}. However, this later interpretation does not hold when $X$ is not finite (and does not behave like a generalized cohomology theory).
\end{recollection}





\begin{remark}[Borel-type theories]
There are also cheaper ways of extending topological $K$-theory to the equivariant context -- one can use the non-equivariant definition on the Borel construction. The resulting theory is simpler, but carries considerably less information. Already for the equivariant point, it returns only the completion of its representation ring at the origin \cite{AS69}.
Also compare to \cite{Tho86} and the discussion \cite[\S 5]{Tho88b}. We will not use the Borel-type version.  
\end{remark}




\subsubsection{Goresky--Kottwitz--McPherson descriptions.}\label{section: GKM descriptions}


Let $G$ be a topological group and $E_G$ a multiplicative $G$-equivariant cohomology theory. Let $X$ be a topological space with $G \acts X$. Under favorable assumptions, the equivariant generalized cohomology $E^{\bullet}_G(X)$ can be computed combinatorially from the fixed points of $G \acts X$ by solving a specific system of congruences. For $T$-equivariant singular cohomology $H^{\bullet}_T(X; \C)$ of a complex algebraic variety $X$, this was realized in \cite[\S 7]{GKM97}. It was then extended to generalized cohomology theories in \cite{HHH05}; we will follow their excellent treatment, to which we refer for details.
\begin{assumption}\label{assumption: HHH assumptions}
Let $G \acts X$ be a topological group acting on a topological space. Assume the following \cite[\S 3, Assumptions 1-4]{HHH05}:
\begin{enumerate}
    \item The space $X$ admits a $G$-equivariant stratification $X = \bigsqcup_{i \in I} X_i$ labeled by a poset $I$, whose graded pieces $X_i / X_{<i}$ are Thom spaces of $G$-equivariant vector bundles $V_i$ over some $G \acts F_i$.
    \item The bundles $V_i$ are $E$-orientable and admit a decompositions $V_i = \bigoplus_{j < i} V_{i, j}$ by $G$-equivariant $E$-orientable vector bundles.
    \item There is a system of $G$-equivariant maps $a_{i, j}: F_i \to F_j$ for each $j < i$, such that the following diagrams commute
    \begin{equation*}
    \begin{tikzcd}
        S(V_{i, j}) \arrow[r] \arrow[d] & X_{<i} \\
        F_i \arrow[r, "a_{i, j}"] & F_j \arrow[u, hook]
    \end{tikzcd}
    \end{equation*}
    In words, the attaching map $S(V_i) \to X_{<i}$ from the sphere bundle of $V_i$ is, on each subbundle $S(V_{i, j})$, induced by $a_{i, j}$.
    \item Fix any $i \in I$. As $j<i$ varies, the Euler classes $e(V_{i, j})$ are not zero divisors and pairwise relatively prime in $E^{\bullet}_G(F_i)$.
\end{enumerate}
\end{assumption}

The main result of \cite{HHH05} then reads as follows.
\begin{proposition}\label{proposition: GKM of HHH}
Assuming (1), the following restriction map is injective
\begin{equation*}
    E^{\bullet}_G(X) \hookrightarrow \prod_{i \in I} E^{\bullet}_G(F_i).
\end{equation*}
Assuming (1)-(4), its image is given by the GKM ring
\begin{equation*}
    \GKM^{\bullet}_G(X) := \big\{ (f_i) \in \prod_{i \in I} E^{\bullet}_G(F_i) \mid f_i - a^*_{i, j} (f_j) \in e(V_{i, j}) E^{\bullet}_G(F_i)  \ \forall j < i \big\}.
\end{equation*}   
\end{proposition}
\begin{proof}
   See \cite[Theorems 3.1 and 2.3]{HHH05}.
\end{proof}

\begin{remark}\label{remark: GKM in the BB case}
The most common use case is that of a torus action on a complex algebraic variety $T \acts X$. It induces a Białynicki-Birula decomposition of $X$ by attractors of fixed points. Assume that the fixed points $X^T$ are isolated and the strata are affine spaces. One then takes $(F_i, i\in I)$ to be the set of fixed-points, $V_{i}$ the attractor cells, $V_{i, j}$ their weight decompositions, $a_{i, j}$ the obvious maps. In many real-life examples -- namely when the $T$-action is sufficiently non-degenerate -- the conditions (3) and (4) hold.
\end{remark}

We have followed \cite{HHH05}. Also see \cite[Appendix A]{Ros03}, \cite{HHRW16}, \cite{Zho25}, \cite{KR21} for related subsequent work.









\subsection{Algebraic cohomology theories as localizing invariants}\label{section: algebraic cohomology theories as localizing invariants}


\subsubsection{Localizing invariants.}
Certain generalized (co)homology theories do have a purely algebraic counterpart as localizing invariants $E(-)$ of stable categories. Given a space $X$ of algebro-geometric origin, one can then compute $E(X)$ via its derived category. More precisely, such theories again come in pairs: the covariant version $E(\Coh(X))$ and the contravariant version $E(\Perf(X))$. 

Again, we are biased towards the latter.
As in the topological case, one good reason for this is the ring structure. Namely, assume $E(-)$ is lax monoidal. Since $\Perf(X)$ has a symmetric monoidal structure via $\otimes$, the output $E(\Perf(X))$ is an $\mathbb{E}_{\infty}$-ring spectrum.

Given an affine group scheme $G$ acting on a qcqs scheme $X$, we put
\begin{equation*}
    E_G(X) := E(\Perf(G \backslash X)),
\end{equation*}
in line with \cite{Tho88, TT90}.
Taking homotopy groups, we obtain the $\Z$-graded abelian group $\pi_{\bullet}(E_G(X))$. If $E(-)$ was lax monoidal, this is a graded-commutative ring.

\begin{convention}\label{remark: regrettable indices}
In this paper, we use the following nonstandard convention. Given a localizing invariant $E$, we denote 
\begin{equation*}
E^{i}(\eX) := \pi_{i}(E(\eX)) := \pi_{i}E(\Perf(\eX))),   
\end{equation*}
in other words, we put the homotopy degree to the upper index. While this is non-standard from the viewpoint of localizing invariants, we are forced to do so in order to match the homology-cohomology conventions from algebraic topology: our invariants are computed from $\Perf(\eX)$ and hence compare to cohomological versions in topology. (The lower-index convention would be given by $\Coh(\eX)$ and correspond to homology.)
\end{convention}



The known localizing invariants either come from vector bundles (e.g. algebraic $K$-theory and its variants) or from differential forms (e.g Hochschild homology $HH$ and its variants). 
%
We write $K(\eX) := K(\Perf(\eX))$ for non-connective algebraic $K$-theory.

Perhaps surprisingly, purely topological invariants of qcqs schemes over $\C$ can be recovered from this machinery. The relevant examples are Blanc's topological $K$-theory $K^{\btop}$ and periodic cyclic homology $HP$.




\subsubsection{Equivariant topological $K$-theory.}
We work over $k=\C$. We denote by $K^{\btop}$ Blanc's non-connective topological $K$-theory \cite{Bla16}. This is a multiplicative localizing invariant of $\C$-linear categories.
In particular, given any affine group $G$ action on a qcqs scheme $X$ over $\C$, we obtain
\begin{equation*}
    K^{\btop}_G(X) := K^{\btop}(\Perf(X/G)).
\end{equation*}
It is an $\mathbb{E}_{\infty}$-ring spectrum.

\begin{remark}
For $X$ qcqs scheme over $\C$, Blanc's topological $K$-theory $K^{\btop}(X)$ agrees with the equivariant topological $K$-theory $K^{\top}(X) \simeq KU(X)$ of the analytification from \S \ref{section: generalized cohomology theories in topology}. More generally, for an action of an algebraic group $G \acts X$ where $X$ is qcqs, it agrees with the genuine equivariant topological $K$-theory $K^{\top}_G(X) = KU_G(X)$ of on the analytification $X(\C)$ in the sense of \S \ref{recollection: genuine topological K-theory}. See \cite{Bla16}, \cite{HLP20}, \cite[Remark 2.27]{Low25} for this identification.

Consequently, we may (and will) identify $K^{\btop}_G(X) \simeq K^{\top}_G(X)$ in such situations, in particular when $X$ is a $G$-equivariant projective variety. However, note that for ind-schemes we use the purely topological $K^{\top}_G$.
\end{remark}


\begin{notation}
We denote by $u \in K_{2}(\pt) = K^{2}(\pt)$ the Bott element in algebraic $K$-theory. 

It maps into topological $K$-theory, giving an an invertible element $u \in K^{\btop, 2}(\pt)$ with inverse $u^{-1} \in K^{\btop, -2}(\pt)$. In particular, multiplication by $u$ is an isomorphism on homotopy groups shifting degree by $2$, making $K^{\btop, \bullet}(X)$ 2-periodic. We write $K^{\btop, *}(X)$ with $* \in \Z/2$ for this $\Z/2$-graded, graded-commutative ring.
\end{notation}



\subsubsection{Equivariant periodic cyclic homology.}
Consider now the localizing invariant $HP(-/k)$ of $k$-linear categories. Following our conventions, we write $HP^G(X / k) := HP(\Perf(X/G) / k)$ for the $G$-equivariant periodic cyclic homology of $X \in \Sch^G_{k}$. 
We write $u \in HP_{2}(\pt / k) = HP^{2}(\pt / k)$ for the periodicity operator with inverse $u^{-1} \in HP^{-2}(\pt)$ as in \cite[\S 5.1.4]{Lod97}.

Let us specialize to $k=\C$. Then the Dennis trace induces a map of localizing invariants $K^{\btop}(-) \to HP(- / \C)$ called the topological Chern character. The trace map $K^{\btop}(\pt; \C) \to HP(\pt / \C)$ identifies the Bott element on the left-hand side with the periodicity operator on the right-hand side. See \cite{Lod97} for classical exposition; see \cite{Chen20, EKS25} for the equivariant case. We will recall its geometric incarnation in terms of fixed points in \S \ref{section: fixed-point schemes} below.


\subsubsection{Equivariant base.}
For any group scheme $G$ over $k$, we tautologically identify the equivariant $K$-theory of a point with the representation ring
\begin{equation*}
    K^0_G(\pt) = K^{\btop, 0}_G(\pt) = R(G)
\end{equation*}
ans similarly with coefficients.
Working over $k = \C$, the topological Chern character identifies
\begin{equation*}
    K^0_G(\pt; \C) = K^{\btop, 0}_G(\pt; \C) = R(G; \C) = HP^0_G(\pt / \C).
\end{equation*}


\subsubsection{Derived fibers over $T$.}
Working with a torus $T$ over $k$, the equivariant base identifies with its ring of functions.
\begin{equation*}
 K^{0}_T(\pt; k) = K^{\btop, 0}_T(\pt; k) = \O(T).   
\end{equation*}
Given any closed point $t \in T$, we denote $t: \O(T)\to k$ the corresponding algebra homomorphism. To emphasize the $\O(T)$-algebra structure on the right-hand side, we sometimes write $k = k_t$. We further denote $(-)_{\lL t} = (-) {\otimes^{\lL}_{\O(T)}} k_t$ the derived tensor product along this map. 

In particular, the functor $(-)_{\lL t}$ can be applied to $K^{\btop}_T(-; k)$ and $HP^{\btop}_T(- / k)$, both regarded as algebra objects in the derived category of $\O(T)$-modules, getting
\begin{equation*}
    K^{\btop}_T(X; k)_{\lL t} \qquad \text{and} \qquad HP_T(X / k)_{\lL t}.
\end{equation*}

\subsubsection{Remarks.}
Let us outline how the current section parallels the topological discussion.

\begin{remark}[Comparisons]
Comparisons of algebraic $K_G$ and topological $K^{\top}_G$ have a long history -- algebraic invariants usually encode the topological ones. Algebraic $K$-theory of coherent sheaves (usually called $G$-theory) relates to topological $K$-homology; algebraic $K$-theory of perfect complexes relates to topological $K$-cohomology. See \cite{Tho88, Tho88b}.
\end{remark} 

\begin{remark}[Kan extended theories]
In the same way as in topology, there is a cheaper way of producing equivariant invariants. Instead of using $\Perf(G \backslash X)$, one can Kan-extend from the non-equivariant case. The result then parallels the Borel-type variant in topology: it is much simpler and sees only the completion at the augmentation ideal. See \cite{Tho86}, or also \cite{TV24}. Again, we do not use these variants.    
\end{remark}


\begin{remark}[Motivic filtrations]\label{remark: motivic filtrations}
Algebraic $K$-theory and its variants of any $X \in \Sch_{k}$ carry the \textit{motivic filtration} \cite{EM23}. In characteristic zero and after taking rational coefficients, this filtration naturally splits. In particular, $K^0(X; \C)$ is naturally a graded ring. There is a compatible filtration on $HP(-/ \C)$. Interpreting the latter in terms of de Rham cohomology, this is precisely the de Rham filtration. Again, it naturally splits, making $HP^0(X/ \C)$ into a graded ring. The trace map is compatible with these filtrations.

Such filtrations also appear in the equivariant setup. For example, one can get a filtration on $K_G(X)$ by Kan extension from the motivic filtration on non-equivariant $K$-theory as in \cite{EKS25}. Also see \cite[(2.6)]{TV24}, \cite{KN25}, \cite{EKS25} for more discussion of Atyiah--Segal completion phenomena in this context. However, it is not well understood whether such a filtration exists genuinely \cite[Question 1.0.2]{EKS25}, in other words ``relatively over $BG$".
\end{remark}










\subsection{Fixed-point schemes}\label{section: fixed-point schemes}


Let $G$ be a group scheme over $k$ and denote $\Sch_k^G$ the category of $G$-equivariant $k$-schemes. We recall the notion of the associated fixed-point scheme from \cite[\S 1]{Low24}.


\subsubsection{Fixed point schemes.}
Suppose $G$ acts on $X \in \Sch^G_k$ via $m: G\times X \to X$.  We define the associated \textit{fixed-point scheme} as the fiber product of schemes
\begin{equation}\label{equation: definition of fixed-point schemes}
\begin{tikzcd}
\Fix_G(X) \arrow[r] \arrow[d] & G\times X \arrow[d, "m \times \pr_2"] \arrow[r, "\pr_1"] & G \\
X \arrow[r, "\Delta_X"] & X\times X
\end{tikzcd}
\end{equation}
Its functor of points is given by
\begin{equation*}
    \Fix_G(X): R \mapsto \{ (g, x) \in G(R) \times X(R) \mid gx = x \}.
\end{equation*}
The horizontal composition in \eqref{equation: definition of fixed-point schemes} gives a map $\pi: \Fix_G(X) \to G$. For any $g \in G(k)$, the fiber $\pi^{-1}(g)$ parametrizes those $x \in X(R)$ which are fixed by $g$. 
Given any $S \in \Sch_k$ with a map $S \to G$, we obtain
\begin{equation}
\Fix_S(X) := S {\times}_G \Fix_G     
\end{equation}
and view it as the fixed point family of the restriction of the $G$-action to $S$. In particular, given any closed point $g \in G$, we have its scheme-theoretic fixed-points $\Fix_g(X) = X^g$.

\subsubsection{Residual action.}\label{section: residual action on Fix}
The fixed-point scheme carries a natural action $G \acts \Fix_G(X)$ induced from the original action on $G \acts X$ and conjugation action $G \acts G$ on the defining diagram  \eqref{equation: definition of fixed-point schemes}. We denote the quotient by $\Fix_{\tfrac{G}{G}}(X)$.

In particular, suppose $h, g \in G$ and denote $g' = hgh^{-1}$. Then we have a natural identification
\begin{equation*}
    \Fix_{g}(X) = \Fix_{hgh^{-1}}(h\cdot X) \xrightarrow{=} \Fix_{g'}(X)
\end{equation*}
where the latter arrow is induced by the translation $h^{-1}: X \to X$.
For fixed $g \in G$, we thus obtain an action
\begin{equation*}
    \Cent_G(g) \acts \Fix_g(X).
\end{equation*}


\subsubsection{Fixed-point schemes preserve abstract blowups.}
We will need the following elementary observation about fixed-point schemes.
\begin{lemma}\label{lemma: fixed-point schemes of abstract blowup squares}
Let $G$ be a group over $k$, and $S \to G$ any map from some $S \in \Sch_k$. Then the classical fixed-point scheme functor
\begin{equation*}
    \Fix_S(-): \Sch_k^G \to \Sch_S
\end{equation*}
preserves abstract blowup squares. In particular, this applies to $\Fix_g(-)$ for any point $g \in G$.
\end{lemma}
\begin{proof}
The classical fixed-point scheme functor $\Fix_S(-)$ preserves isomorphisms, closed immersions, open immersions, closed-open decompositions, proper maps \cite[Propositions 1.2 and 1.3]{Low24}. It commutes with fiber products (in fact with all limits \cite[Lemma 1.7]{Low24}). Therefore, it also preserves abstract blowup squares.
\end{proof}



\subsubsection{Reduced and derived variants.}\label{subsubsection: reduced and derived variants}
If $\Fix_\frac{G}{G}(X)$ is non-reduced, we can take its reduction $\Fix^{\red}_{\frac{G}{G}}(X)$. On the other hand, we may take the defining fiber product \eqref{equation: definition of fixed-point schemes} in derived schemes, giving a derived enhancement $\Fix^{\lL}_{\frac{G}{G}}(X)$. We get closed immersions
\begin{equation}\label{equation: reduced, classical and derived fixed-point schemes}
    \Fix^{\red}_{\frac{G}{G}}(X) \hookrightarrow \Fix_{\frac{G}{G}}(X) \hookrightarrow \Fix^{\lL}_{\frac{G}{G}}(X).
\end{equation}
For the applications in this paper, the reduced version will be sufficient.


On the other hand, the derived variant has the following interpretation. 
For any derived stack $\eX$ over $k$, its derived loop space $\dL (\eX)$ is defined by the derived self-intersection of the diagonal $\dL (\eX) = \eX \times_{\eX \times_k \eX} \eX = \Map(S^1, \eX)$. Applied to $\eX = G \backslash \eX$, there is an equivalence
\begin{equation*}
    \dL (G \backslash X) \simeq \Fix^{\lL}_{\frac{G}{G}}(X).
\end{equation*}
This interpretation allows to compare to periodic cyclic homology, see Proposition \ref{proposition: periodic cyclic homology and derived fixed-point schemes} below. If $G$ is an affine group acting on a quasi-projective variety $X$, the derived structure on $\Fix^{\lL}_{\frac{G}{G}}(X)$ is trivial if and only if $G\acts X$ acts with finitely many orbits by the criterion \cite[Proposition 3.3]{BN13}.








\subsection{\texorpdfstring{$K$}{K}-theory, periodic cyclic homology, and cohomology of fixed points}\label{section: K-theory and HP and cohomology of fixed points}

\subsubsection{Localization theorem.}
We work over the base field $k=\C$. Let $T = \Gm^{n}$ be an algebraic torus over $k$. 
Throughout this section, we write $\Sch^T_{\C} := \Sch^{T, \qcqs}_{\C}$ for the category of $T$-equivariant, quasi-compact quasi-separated schemes over $\C$  -- we omit the superscript $\qcqs$ for readability. 


Let $X \in \Sch^{T}_{\C}$ be any $T$-equivariant qcqs scheme over $\C$; it is in particular allowed to be singular. 
The purpose of this section is to recall the (maybe not so well) known relationship between its complexified equivariant topological $K$-theory $K^{\btop}_T(X; \C)$, equivariant periodic cyclic homology $HP^T(X / \C)$ and the singular cohomologies of the varying fixed-point subschemes $\Fix_t(X)$ for $t \in T$.

\begin{theorem}\label{theorem: fibers of topological K-theory, HP and fixed-points}
Let $X \in \Sch^{T}_{\C}$. Then we have equivalences
\begin{equation}\label{equation: K, HP, Fix}
    K^{\btop}_{T}(X; \C) \xrightarrow{\simeq} HP_T(X / \C) \xrightarrow{\simeq} \RG(\Fix^{\lL}_{\frac{T}{T}}(X), \O)^{tS^1}
\end{equation}
in $\CAlg(\O(T)[u^{\pm 1}])$. 
%
Taking fibers over a closed point $t \in T$, this equivalence specializes to
\begin{equation}\label{equation: fibers and singular cohomology of fixed-points}
    K^{\btop}_T(X; \C)_{\lL t} \xrightarrow{\simeq} H_{\sing}(\Fix_t(X); \C)\llp u \rrp
\end{equation}
in $\CAlg(D(\C))$.
\end{theorem}
\begin{proof}
Proposition \ref{theorem: fibers of HP and cohomlogy of fixed points}, Proposition \ref{proposition: periodic cyclic homology and derived fixed-point schemes} and Proposition \ref{theorem: fibers of equivariant K-theory nd cohomology of fixed points} below.
\end{proof}


\begin{remark}
In plain english: the fibers of equivariant topological $K$-theory over various closed points $t$ of the equivariant base $T$ are isomorphic with the singular cohomologies $H_{\sing}(\Fix_t(X); \C)\llp u \rrp$ of the corresponding scheme-theoretical fixed-points $\Fix_t(X)$. Reformulating, $T$-equivariant topological $K$-theory of $X$ puts the singular cohomologies of the fixed-point varieties $\Fix_t(X)$ into an algebraic family.
This algebraic family is geometrically computed by 
$$\RG(\Fix^{\lL}_{\frac{T}{T}}(X), \O))^{tS^1},$$
the Tate fixed points of the circle action on global functions on the derived fixed-point stack $\Fix^{\lL}_{\frac{T}{T}}(X)$.
\end{remark}

We also have the following formally deformed variant on the final assertion of the above theorem.
\begin{proposition}\label{remark: the completed variant of equivariant localization}
Let $X \in \Sch^{T}_{\C}$. Taking completions $(-)^{\wedge}_{\lL t}$ along $t \in T$, we obtain
\begin{equation}\label{equation: completed version K-theory}
  K^{\btop}_T(X; \C)^{\wedge}_{\lL t} \xrightarrow{\simeq} H_{T}(\Fix_t(X); \C) \llp u \rrp^{\wedge} 
\end{equation}
Here, the completion of the right-hand side it with respect to the derived Hodge filtration in the Cartan model of equivariant cohomology. 
\end{proposition}
\begin{proof}
Take Proposition \ref{theorem: topological K-theory and HP}, then take the formal derived completion $(-)^{\wedge}_{\lL t}$ over $t \in T$, and finally apply Proposition \ref{proposition: fibers of HP and singular cohomology -- deformed}.
\end{proof}



\subsubsection{Case of reductive groups.}

\begin{remark}[Case of reductive groups.]
Let $G$ be any reductive group acting on a quasi-projective variety $X$ over $\C$. Then the equation \eqref{equation: K, HP, Fix} holds $G$-equivariantly:
\begin{equation*}
    K^{\btop}_{G}(X; \C) \xrightarrow{\simeq} HP^G(X / \C) \xrightarrow{\simeq} \RG(\Fix^{\lL}_{\frac{G}{G}}(X), \O)^{tS^1}.
\end{equation*}
Moreover, let $g \in G$ be an element with semisimple part $\overline{g} \in \Spec(R_{\C}(G))$. Then the identification \eqref{equation: fibers of HP and singular cohomology of fixed-points} goes through as:
\begin{equation*}\label{equation: fibers and singular cohomology of fixed-points semisimplification}
    K^{\btop}_T(X; \C)_{\lL \overline{g}} \xrightarrow{\simeq} H_{\sing}(\Fix_{\overline{g}}(X); \C)\llp u \rrp
\end{equation*}
Note that this does {not} work without semisimplicifation.
\end{remark}
\begin{proof}[Proof of Remark.]    
This version follows by the same arguments, combining results of \cite{Chen20} in the smooth case with the recent progress on equivariant cdh descent \cite{LS25} with respect to reductive groups $G$. Also see \cite[Appendix B]{Low25}.
\end{proof}

In particular, even if we consider $G$-equivariant $K$-theory instead of $T$-equivariant one, there is no extra information in its fibers.


\subsubsection{$K$-theory and periodic cyclic homology.}
Let us first recall the relationship between equivariant topological $K$-theory and equivariant periodic cyclic homology. 

\begin{proposition}\label{theorem: topological K-theory and HP}
For any $X \in \Sch^{T}_{\C}$ we have an equivalence
\begin{equation*}
K^{\btop}_T(X; \C) \xrightarrow{\simeq} HP_T(X / \C)    
\end{equation*}
as algebras over $K_T^{\btop}(\pt; \C) \simeq HP_T(\pt / \C)$.
\end{proposition}
\begin{proof}
For a sufficient reference covering (possibly singular) stacks with nice stabilizers, see \cite{Kha23}. This heavily builds on the previous works \cite{Bla16, Kon21, HLP20, ES21}. 
\end{proof}





\subsubsection{Periodic cyclic homology and fixed-points.}
Periodic cyclic homology $HP(-/\C)$ over $\C$ is further geometrically modeled by suitable functions on the derived fixed-point scheme.
\begin{proposition}\label{proposition: periodic cyclic homology and derived fixed-point schemes}
There is a natural equivalence
    \begin{equation*}
        HP_T(X/\C) \xrightarrow{\simeq} \RG(\Fix^{\lL}_{\frac{T}{T}}(X), \O)^{tS^1}.
    \end{equation*}
Here, $(-)^{tS^1}$ refers to the Tate construction with respect to the natural $S^1$-action on the derived fixed-point scheme.
\end{proposition}
\begin{proof}
Apply the Tate construction to \cite[Discussion 1.12]{Low24} and \cite[Example 2.2.20]{Chen20}
\end{proof}







\subsubsection{Periodic cyclic homology computes cohomology of fixed-points.}
Let us now explain how fibers of equivariant periodic cyclic homology relate to the singular cohomology of fixed-points.
\begin{proposition}\label{theorem: fibers of HP and cohomlogy of fixed points}
Let $t \in T$ be a closed point. We then have an equivalence
\begin{equation}\label{equation: fibers of HP and singular cohomology of fixed-points}
    HP_T(X/ \C)_{\lL t} \xrightarrow{\simeq} H_{\bet}(\Fix_t(X); \C)\llp u \rrp.
\end{equation}
\end{proposition}
\begin{proof}
First recall the identification of derived de Rham cohomology with singular cohomology of the analytification \cite[Theorem 4.2.9 and Corollary 4.2.11]{Chen20} and \cite[Theorem 4.10]{Bha12}, valid for any derived stack of finite type over $\C$ (the case of quasi-projective varieties is sufficient for us).

We first assume that the $T$-variety $X$ is smooth. Then \eqref{equation: fibers of HP and singular cohomology of fixed-points} holds by \cite[Corollary 1.0.4]{Chen20} (also see \cite[Theorem D]{Chen20} for a finer statement), using the above discussion to identify the right-hand side with singular cohomology. We use here the semisimplicity assumption on $t$.

Now, both sides of \eqref{equation: fibers of HP and singular cohomology of fixed-points} satisfy proper excision on $T$-equivariant abstract blowup squares.
Indeed, the left-hand side is given by the composition
\begin{equation*}
    \Sch^{T}_{\C} \xrightarrow{HP_T(-/ \C)} \D(HP_T(\pt / \C)) \xrightarrow{(-)_{\lL t}} \D(\C)
\end{equation*}
Here, $HP(-/\C)$ is a truncating invariant, so $HP^T(-/\C)$ has proper excision by \cite{ES21}. The same holds after composition with the derived functor $(-)_{\lL t} = (-)\otimes^{\lL}_{HP_T(\pt / \C)} \C_t$.
On the other hand, the right-hand side can be factored as
\begin{equation*}
    \Sch^{T}_{\C} \xrightarrow{\Fix_t(-)} \Sch_{\C} \xrightarrow{H_{\bet}(-; \C)\llp u \rrp} D(\C).
\end{equation*}
Here, the functor $\Fix_t(-)$ preserves abstract blowup squares by Lemma \ref{lemma: fixed-point schemes of abstract blowup squares}. 
Moreover, singular cohomology (with any coefficients) has proper excision (this being the case for any generalized cohomology theory with open-closed long exact sequence). Alternatively, this last claim can be seen on the level of derived de Rham cohomology, which is a canonical direct summand of $HP(-/\C)$ and hence has cdh descent \cite[Lemma 4.5]{EM23}.

Altogether, by equivariant resolution of singularities in characteristic zero and induction on dimension, the singular case of \eqref{equation: fibers of HP and singular cohomology of fixed-points} follows from the smooth one.
\end{proof}

\subsubsection{Completions of equivariant cohomology.}
The above proposition has the following formally deformed variant.
\begin{proposition}\label{proposition: fibers of HP and singular cohomology -- deformed}
Let $X \in \Sch^{T}_{\C}$. For any closed point $t \in T$, we have
\begin{equation}\label{equation: completed version HP}
  HP_T(X; \C)^{\wedge}_{\lL t} \xrightarrow{\simeq} H_{T, \sing}(\Fix_t(X); \C) \llp u \rrp^{\wedge}.
\end{equation}
\end{proposition}
\begin{remark}
Here, the completion of the left-hand side is the derived formal completion over $t \in T$. The completion of the right-hand side is with respect to the derived Hodge filtration in the Cartan model of equivariant cohomology -- see \cite[Example 2.1.10, Definition 4.2.3]{Chen20} for a discussion (using the cotangent complex in general).
\end{remark}
\begin{proof}
If $X$ is smooth, the statement is (a special case of) \cite[Theorem D, Theorem 4.3.2]{Chen20}. We again deduce the general case by proper excision. For the left-hand side, $HP(-/\C)^{\wedge}_{\lL t}$ satisfies proper excision -- it is given by composing the truncating invariant $HP(-/ \C)$ with $(-)^{\wedge}_{\lL t}$. For the right-hand side, we again use that $\Fix_t(-)$ preserves abstract blowups. Moreover, since the derived de Rham complex satisfies proper excision, the associated filtrations on $H_{\sing}(-; \C) \llp u \rrp$ are also compatible with proper excision -- so taking $H_{\sing}(-; \C)\llp u \rrp^{\wedge}$ satisfies proper excision. By existence of equivariant resolutions in characteristic zero, we thus reduce to the smooth case.
\end{proof}


\subsubsection{Fibers of equivariant $K$-theory and cohomology of fixed-points.}
Putting the above discussion together, we obtain the second part of the theorem.
\begin{proposition}\label{theorem: fibers of equivariant K-theory nd cohomology of fixed points}
    Let $X \in \Sch^T_{\C}$ and $t \in T$ be a closed point. Then we have an equivalence
\begin{equation*}
    K_T^{\btop}(X; \C)_{\lL t} \xrightarrow{\simeq} H_{\sing}(\Fix_t(X); \C)\llp u \rrp
\end{equation*}
in $\CAlg(D(\C))$.
\end{proposition}
\begin{proof}
Let $t: \C[t_1^{\pm 1}, \dots, t_n^{\pm 1}] \to \C$ be a classical point of the equivariant base. Taking the fiber $(-)_t$ of the equivalence from Proposition \ref{theorem: topological K-theory and HP} and composing with the equivalence from Proposition \ref{theorem: fibers of HP and cohomlogy of fixed points}, we deduce
\begin{equation*}
K_T^{\btop}(X; \C)_{\lL t} \xrightarrow{\simeq} HP_T(X / \C)_{\lL t} \xrightarrow{\simeq} H_{\bet}(\Fix_t(X); \C)\llp u \rrp.   
\end{equation*}
\end{proof}

\subsubsection{Equivariant formality.}
Under extra flatness assumptions, the above derived identifications simplify. In particular, this happens under suitable equivariant formality hypothesis. Let us hide the formal parameter $u$ by regarding $K_{T}^{\btop, *}(X)$ as graded by $* \in \Z/2$.

\begin{corollary}\label{corollary: K-theory in the flat situation}
Assume that $K_{T}^{\btop, *}(X)$ is a flat module over the equivariant base and let $t \in T$.
Then we have an isomorphism of $\Z/2$-graded-commutative rings
\begin{equation*}
    K_{T}^{\btop, *}(X; \C) \underset{\O(T)}{\otimes} \C_t = H^{\bullet}_{\bet}(\Fix_t(X); \C)
\end{equation*}
where $* \in \Z/2$ and $\bullet \in \N_{\geq 0}$. 
\end{corollary}
\begin{proof}
Take $\pi_{\bullet}(-)$ in Proposition \ref{theorem: fibers of equivariant K-theory nd cohomology of fixed points}.
By the flatness assumption, the base change along $\O(T) \to \C_t$ is underived, hence commutes with the passage to homotopy groups on the left-hand side. 

Both sides are $2$-periodic; the trace map is compatible with the natural filtrations; the element $u$ has degree $-2$. This explicitly identifies
\begin{equation*}
    K_{T}^{\btop, 0}(X; \C) \underset{\O(T)}{\otimes} \C_t = \bigoplus_{i \geq 0} H^{2i}_{\sing}(\Fix_t(X); \C) \cdot u^i
\end{equation*}
\begin{equation*}
    K_{T}^{\btop, 1}(X; \C) \underset{\O(T)}{\otimes} \C_t = \bigoplus_{i\geq 0} H^{2i+1}_{\sing}(\Fix_t(X); \C) \cdot u^i
\end{equation*}
\end{proof}

\begin{remark}\label{remark: K-theory in the superflat situaton}
Even more specifically, suppose that $K_{T}^{\btop, *}(X)$ is concentrated in the even part, where it is given by a flat module over the equivariant base. Then we have a ring isomorphism
$$K_{T}^{\btop, 0}(X; \C)\underset{\O(T)}{\otimes} \C_t =  H^{\bullet}_{\bet}(\Fix_t(X); \C),$$
the right-hand size is concentrated in even degrees, and the motivic grading matches the cohomological grading. This is the case relevant for our applications.
\end{remark}


\begin{aside}
An appropriate definition of the equivariant motivic filtration (Remark \ref{remark: motivic filtrations}) on the left-hand side would match the cohomological filtration on the right-hand side. For example, this is the case over $1 \in T$ via the Kan-extended motivic filtration.
\end{aside}



\subsubsection{Bibliographical remarks.}
The groundwork for the presented statements is contained in the literature. Moreover, weaker versions of some of these claims can be recovered from the classical literature on completion and localization techniques in equivariant topology, pioneered by Segal \cite{Seg68} and Goresky--Kottwitz--MacPherson \cite{GKM97}. These weaker statements would be also sufficient for our main applications. Nevertheless, we hope our exposition gives a conceptually clear and useful viewpoint. 

Our discussion above is based on the following two lines of development.
\begin{itemize}
    \item The relationship between topological $K$-theory $K^{\btop}$ and periodic cyclic homology $HP$ comes, in various degrees of generality, from \cite{Bla16}, \cite{HLP20}, \cite{Kon21}, \cite{ES21}, \cite{Kha23}, \cite{LS25}. 
    \item The relationship between $HH$ and functions on derived fixed-point schemes originates in \cite{BFN10}; the topological version with $HP$ follows by taking the Tate construction. We refer to \cite{Chen20} for a throughout and highly recommendable discussion of equivariant $HP$ with its relation to fixed-point schemes and ultimately to singular cohomology of fixed-points. Also see \cite{Bha12}, \cite{ACH19}, \cite{HSS17}, \cite{BN12, BN13}, \cite{Toe14} for more context.
\end{itemize}
An additive analogue of our fixed-point schemes together with its relation to equivariant cohomology was studied in \cite{HR23} on examples with trivial derived structure.